\section{Introduction}
\label{sec:intro}

\begin{figure*}[t]
  \centering
  \includegraphics[width=0.88\textwidth]{teaser_triptych.png}
  \caption{A facelift-style jaw--neck edit from one photograph (GPT Image~2).
  The prompt-only output (center) changes regions throughout the image; the
  masked composite (right) retains the model output only inside the blend mask.}
  \label{fig:teaser}
\end{figure*}

Surgeons have used simulated outcomes in patient consultations for decades.
Computer imaging has been associated with improved rhinoplasty patient
satisfaction~\cite{sharp2002computerimaging}, and how closely a 3D simulation
matches the eventual outcome predicts how satisfied the patient will
be~\cite{yamamichi2025crisalix}. The same literature warns that an over-idealized
preview sets expectations no operation can meet~\cite{agarwal2007morph}. Producing
these previews today requires either manual morphing skill or dedicated 3D
systems; meanwhile, general-purpose image-editing APIs, i.e., commercial editing
models reachable only through a public endpoint, can perform instruction-guided
photo edits. The natural question is whether such a service can produce a
controlled cosmetic-surgery preview when neither its weights nor its inference
process is accessible.

In preliminary tests, a request for a specific visible change, such as jawline
and neck repositioning or dorsal-hump reduction, also altered skin, lighting, and
other facial regions. Off-target change of this kind can make the output less
recognizable as the same person, and it shows the patient cosmetic effects
unrelated to the requested procedure. The
training-free localization literature knows how to confine
edits~\cite{avrahami2022blended,couairon2023diffedit}. Every such method, however,
needs the model's latents, attention, or sampling loop, and a public editing
API exposes none of them.

We therefore evaluate how much off-target change can be prevented on the
client, i.e., the caller's side of the API, before a request is sent or after
the response returns.
The pilot benchmark covers six editing configurations (GPT Image~2 at two quality
tiers, Nano Banana~Pro and~2, Seedream~5.0 Lite, and FLUX.2~[pro]) plus a Qwen
inpainting endpoint, two requested edit types, and publicly accessible
before/after photograph pairs; each matrix cell contains one generated output.
Our contributions are:

\begin{itemize}
  \item \textbf{A model-agnostic control ladder.} Three rungs over the same
  black-box endpoints: prompt-only, a client-side masked composite that cuts
  the model's edit out along a landmark-derived region mask and pastes it back
  onto the original photograph (with alignment, feathered blending, and pose
  gating), and native masked inpainting where supported. The composite adapts the final
  step of blended-diffusion editing to a setting where nothing inside the model
  is reachable (Sec.~\ref{sec:method}).
  \item \textbf{A two-part diagnostic protocol.} ArcFace identity is paired with
  a CIELAB region-localization ratio and its target/off-target components. The
  metrics detect identity loss, off-target pixel change, and near-copies, but do
  not measure clinical plausibility. Postoperative photographs provide an
  identity reference with a per-face input baseline (Sec.~\ref{sec:protocol}).
  \item \textbf{A pilot benchmark.} Client-side compositing enforces off-mask
  preservation across the tested editors without materially reducing on-target
  pixel change. The editors differ in edit strength and identity retention, while
  the one tested mask-specific endpoint does not outperform compositing. A
  two-order, single-face chain is reported only as a feasibility probe
  (Sec.~\ref{sec:results}).
\end{itemize}

To our knowledge, this is the first controlled comparison of client-side region
confinement across multiple commercial, black-box editing APIs for
cosmetic-surgery previews. What differs here is the setting. Bespoke surgical
simulators and composited rhinoplasty systems
exist~\cite{knoedler2024rhinoplastygan,agarwal2026envisage}, but they assume a
model the operator hosts, and image-editing APIs have been benchmarked in other
domains~\cite{qian2025giebench,gptimgeval2025} without control as the
manipulated variable. When the editor is a black-box endpoint, no
internals-based localization method applies, so whatever control exists must
live on the client. Measuring how much that recovers is our contribution; we
propose no new compositing algorithm, clinical outcome model, or clinical
validation study (Sec.~\ref{sec:related}).\footnote{Analysis code and paper source:
\url{https://github.com/suxrobGM/localize-dont-beautify}}
